\section{Localization and compactness of the curves}
\label{sec:localization}

The rectangle estimates must be transferred to the specified half-plane
measure before they can control its loops.  We first couple matching
states in a bounded window, uniformly in the lattice mesh.  This coupling
does not determine connections outside the window.  Nevertheless, it
preserves the local directed paths needed for both estimates.  We then
use the traversal bound to obtain compactness in the fixed disk
coordinates, retaining every nonconstant limiting loop.

\subsection{A uniform comparison with finite rectangles}

Write $\mathcal R_{\delta,R}$ for a Temperleyan rectangle whose coarse
bottom side is at height $\delta$, whose other sides are within $O(\delta)$
of $x=-R$, $x=R$, and $y=R$, and whose lower left coarse corner is deleted.
Rounding the sides to the coarse grid gives odd numbers of fine-grid
vertices in both directions.  The precise rounding will be immaterial.
Agreement of matching states on a set means agreement on every edge
having at least one endpoint in that set.

\begin{proposition}[Uniform localization]
\label{prop:localization}
Let $K$ be a bounded closed subset of $\overline{\HH}$ and let $\eps>0$.
There is an $R<\infty$ such that, for every sufficiently small $\delta$,
the two independent half-plane matchings can be coupled with two
independent uniform matchings of $\mathcal R_{\delta,R}$ so that both
matching states agree on $K$ with probability at least $1-\eps$.
\end{proposition}

Here is the finite graph to which we apply the Temperley
correspondence~\cite[Theorem~1]{KPW}. Let $G$ be the coarse rectangular
grid of mesh $2\delta$, including its lower left corner $r$, and let
$G^*$ be its plane dual, with exterior vertex $f_\infty$. The fine-grid
vertices are the vertices of $G$, its edge midpoints, and its bounded
face centers. A midpoint is joined to the endpoints of its coarse edge
and to its incident bounded face centers. Deleting $r$ gives exactly
$\mathcal R_{\delta,R}$; the exterior face-node is omitted from this
matching graph.

A spanning tree $T$ of $G$ is directed toward $r$. Its complementary
dual tree $T^*$ is directed toward $f_\infty$. Each remaining primal
vertex is matched to the midpoint of its outgoing edge in $T$.
Each bounded face center is matched to the midpoint of the primal
edge crossed by its outgoing edge in $T^*$. The corner $r$ is incident
with the exterior face, so the
correspondence is a bijection. All edge weights are one: uniform
matchings correspond to uniform primal trees and, by complementation,
to uniform dual trees. In $G^*$ we retain a separate edge to
$f_\infty$ for every primal boundary edge, including parallel edges.

Wilson's algorithm~\cite{Wilson,PW} samples this finite dual tree by
loop-erasing walks stopped on the tree already drawn, starting with
$f_\infty$. Each walk chooses uniformly among the four edge options at
a bounded face. A proposed step crossing a primal boundary edge is
absorbed at $f_\infty$, and we retain that boundary edge as its physical
contact. Theorems~13 and~16 of~\cite{PW} permit any fixed order of
starting sites. We will confine the dual paths needed near $K$ and
then recover the primal directions from the cycles they determine.
The following walk estimate supplies the uniform attachment
probability used in that confinement.

\begin{lemma}[Attachment near an explored branch]
\label{lem:walk-attachment}
Fix $1\le C_0<\infty$.  There are constants $C_1>C_0$ and $p>0$ with the
following property.  Let a square-lattice walk have mesh $a_0$, start at
$x$ above a horizontal killing boundary, and let $a\ge a_0$.
Suppose an already explored nearest-neighbor path starts within distance
$C_0a$ of $x$ and runs to that boundary.  The probability that the walk
hits the path or the killing boundary before leaving $B(x,C_1a)$ is at
least $p$.  The constants are independent of the path, the mesh, and
the distance of $x$ from the boundary.
\end{lemma}

\begin{proof}
Adjoin the closed lower half-plane to the explored path.  The resulting
connected set meets $B(x,C_0a)$ and extends outside every larger ball.
It is therefore enough to give the unrestricted walk a fixed positive
probability of tracing a closed circuit surrounding $B(x,C_0a)$ before
leaving $B(x,C_1a)$.  If such a circuit crosses the killing boundary,
the killed walk has already stopped, which is also a success.

Here is a concrete way to force an actual circuit, rather than just a
near-return.  At unit scale, first cross a small rectangle to the right
of the inner ball horizontally.  Follow a thin polygonal corridor once
around the inner ball and return to cross that same rectangle vertically.
The two crossings meet by planarity; the intervening part of the trace
contains a closed curve of nonzero winding around the inner ball.  All
corridors lie in a fixed larger ball.  The event can be specified by
finitely many crossings with fixed positive margins.  Brownian motion
has positive probability of following these corridors, and the
invariance principle gives a uniform lower bound when $a/a_0$ is large.
Here the polygonal square walk has mean-zero increments with covariance
$\tfrac12 I$. Scalar Donsker convergence
\cite[Theorems~8.1.4 and~8.1.11]{Durrett} gives tightness of its two
coordinate processes; the vector central limit theorem for disjoint
increment blocks identifies their joint limit as planar Brownian
motion with covariance $\tfrac12 tI$. The prescribed corridor event
contains a uniform open neighborhood of a continuous polygonal route
with these crossings and positive margins. Its Brownian probability is
positive, so the open-set Portmanteau bound yields the claimed lower
bound. This comparison does not depend on the absorbing path.
For bounded $a/a_0$, force an appropriate finite square-lattice circuit;
after enlarging $C_1$, only boundedly many steps and finitely many local
lattice configurations are involved.  Decreasing the lower bound to
cover these cases gives $p>0$ for all scales.  A separating circuit must
meet the connected absorbing set, proving the assertion.
\end{proof}

\begin{proof}[Proof of Proposition~\ref{prop:localization}]
We first compare one matching in $\mathcal R_{\delta,R}$ with one in a
larger rectangle.  Choose a fixed rectangular window $W_\infty$ containing
$K$ and a neighborhood of it. For every coarse edge whose midpoint is
incident to a fine-grid edge having an endpoint in $K$, explore the
dual paths from all bounded faces incident with that coarse edge.
These are the finite endpoints of its dual edge; they lie within
$O(\delta)$ of $K$ and hence in $W_\infty$ for small $\delta$.
The explored paths determine the relevant dual edges and, through the
fundamental cycles used below, the relevant primal directions.
The walk has mesh $2\delta$ until it is absorbed. At the bottom we
record the crossed primal edge, using the boundary-contact convention
above.

\paragraph{Nets and shrinking windows.}
We explore coarse nets before finer ones. Each finer starting site
will then be close to a previously drawn path to the bottom. If the
net spacing is $a_j$ and the allowed excursion is $b_j$, a walk must
miss order $b_j/a_j$ opportunities to attach before making such an
excursion. We choose these scales so that the failure probabilities
remain summable after counting all net sites.
Put $a_j=2^{-j}$ and $b_j=B2^{-j/2}$, where $B>0$ is fixed.  Enlarge
$W_\infty$ horizontally and upward by $\sum_{\ell>j}b_\ell$ to obtain
$W_j$.  Thus $W_{j-1}$ has a margin $b_j$ around $W_j$, apart from the
common bottom boundary.  In $W_j$ choose a dual-lattice net $N_j$ of
spacing comparable to $a_j$, including sites next to the bottom.  It
has at most $Ca_j^{-2}$ sites and covers $W_j$ within distance $Ca_j$.
Run Wilson's algorithm from these nets in increasing order, starting
with $j=0$ and skipping sites already in the tree.  Stop at the unique
level $J=J(\delta)$ for which $2\delta\le a_J<4\delta$, and take every
dual site in $W_J$.  Its cardinality is still $O(4^J)$.  Thus every
required site has been explored.

Suppose all branches already drawn end on the bottom.  At a late level
$j$, declare an error if a new walk travels distance $b_j/3$ from its
start before attaching.  Before this happens, the walk stays inside
$W_{j-1}$.  At any intermediate location there is a site of $N_{j-1}$
within $Ca_{j-1}=O(a_j)$, together with its previously drawn path to the
bottom.  Lemma~\ref{lem:walk-attachment} gives a uniformly positive
chance of attachment before the walk leaves a ball of radius $Da_j$.
Before the declared error, the walk is at least $2b_j/3$ from the
other three sides of $W_{j-1}$. Since $Da_j/b_j\to0$, these test balls
stay within those sides for all sufficiently large $j$; they may
cross the killing boundary.
By the strong Markov property at successive ball exits,
traveling distance $b_j/3$ without attachment requires at least
$cb_j/a_j$ failed trials.  Consequently its conditional probability is
at most
\begin{equation}
 C\exp(-c b_j/a_j)\le C\exp(-c2^{j/2}).
 \label{eq:wilson-net-tail}
\end{equation}
Independence between these trials or between different walks is not
needed.  The conditional lower bound on attachment suffices.

Summing over all starts at levels above $j_0$ bounds the late-level
error by
\begin{equation}
 C\sum_{j>j_0}4^j\exp(-c2^{j/2}),
 \label{eq:wilson-net-sum}
\end{equation}
uniformly in $\delta$.  This tends to zero as $j_0\to\infty$.
The final all-sites level obeys the same estimate.

\paragraph{The finitely many early walks.}
Fix $j_0$ so that \eqref{eq:wilson-net-sum} is small.  There are boundedly
many starts up to that level, all in a fixed bounded window.  An
unrestricted square-lattice walk started at height at most $H$ obeys,
in diffusive time units,
\[
 \PP(\tau_{\mathrm{bottom}}>T)
 \le \frac{C(H+\delta)}{\sqrt T},
 \qquad
 \PP\left(\max_{s\le T}|X_s-X_0|>A\right)
 \le \frac{CT}{A^2}.
\]
The first inequality is the one-dimensional hitting estimate for the
lazy vertical coordinate; the second is the martingale maximal
inequality.  Choosing first $T$ and then $A$ makes the probability of
leaving a large window before hitting the bottom arbitrarily small,
uniformly for fine meshes.  A Wilson walk stops no later if it first
hits an earlier branch.  A union bound therefore confines all early
walks with arbitrarily high probability.

Choose $R$ to contain this large window and the net windows with their
allowed excursions.  On the event of no early or late error, every
required dual branch lies in $\mathcal R_{\delta,R}$ and reaches the
wired exterior through its bottom side.  Use the same walk increments
in $\mathcal R_{\delta,R}$ and any larger exhausting rectangle.
Their explorations agree exactly on this event.  Adding the remaining
Wilson branches does not alter paths already drawn.

\paragraph{Primal directions.}
It remains to recover more than dual edge membership.  The matching
also uses the directions of primal tree edges toward the deleted
corner.  Let $e$ be a primal tree edge relevant to $K$.  Its crossing
dual edge $e^*$ is absent from the dual tree.  Adding $e^*$ creates a
dual fundamental cycle, which separates the two components of the
primal tree with $e$ removed.  The direction of $e$ is from the
component not containing the root to the component containing it.
The dual branches from the endpoints of $e^*$ determine this cycle;
these were among the required explored branches.

If the cycle avoids the wired exterior, it is contained in the bounded
exploration window, and the remote corner root lies outside it.  If it
uses the wired exterior, all its exterior contacts are on the bottom
within that same bounded window.  Represent the exterior part just
below the bottom, between the two contacts.  The cycle cuts off a
bounded pocket along that bottom interval; the remote corner root is
outside the pocket.  A bottom primal edge is treated identically, with
one contact supplied by $e^*$ itself.  No contact on a remote side or
top can occur on the confinement event.  Thus the component containing
the root, and hence every required primal direction, is the same in
the two rectangles.  Planar duality and the Temperley correspondence
now give identical matching states on $K$.

For each fixed mesh the number of matching states on $K$ is finite.
Let the comparison rectangle exhaust $\HH$ and take a subsequential
limit of the coupled window laws.  Its second marginal is the specified
half-plane local limit, and the agreement bound persists.  Extending
this window coupling by the respective conditional laws gives the
coupling of full matchings.  Finally use two independent copies of the
construction, with half the allotted error for each copy.  This
preserves independence within each pair and proves the proposition.
\end{proof}

\subsection{The two geometric consequences in the half-plane}

We next state exactly what is preserved when the connections outside a
window change.  A directed loop always has the orientation fixed by the
matching labels: first-matching edges go from black to white and
second-matching edges from white to black.  For a directed path portion
$\alpha$ and a transverse segment $I$, retain the notation $\iota_I(\alpha)$ for its
signed intersection sum.  The path endpoints are kept off the test
line, and the segment endpoints off the lattice graph.  A fixed line
between lattice columns or rows may be used as a seam; it need not
bisect the crossed edges.

\begin{corollary}[Local traversal counts]
\label{cor:local-crossings}
Fix a bounded window $K\subset\overline{\HH}$ and two distinct parallel
coordinate lines.  For horizontal lines assume both have positive
height.  The maximum total number of parameter-disjoint portions of
half-plane double-dimer cycles that lie in $K$ and traverse from one
line to the other, in either direction, is tight as $\delta\downarrow0$.
\end{corollary}

\begin{proof}
Apply Proposition~\ref{prop:localization} to a fixed neighborhood of
$K$.  On its agreement event, all such portions are also portions of
rectangle cycles.  Their connections
outside $K$ may differ, but their crossing paths and their
parameter-disjointness do not.  Their number is therefore bounded by
the rectangle traversal count of
Proposition~\ref{prop:rectangle-crossings}.  Take $R$ large enough that
both cuts have positive longitudinal room to the rectangle ends.
For horizontal cuts this uses their positive heights; for vertical
cuts it follows by increasing $R$.  The coupling error is arbitrary,
so the rectangle tightness gives the assertion.  If an original cut
contains lattice vertices, use two slightly inward seams instead:
every traversal of the original slab crosses the inward seams, whose
limiting separation is unchanged.
\end{proof}

\begin{proposition}[Local exclusion of opposite intersection signs]
\label{prop:local-opposite}
Let $u$ be either the horizontal or vertical coordinate, let
$L=\{u=c\}$ be a test line, and let $J\subset L$ be a bounded segment.
If $u$ is the vertical coordinate, assume $c>0$.  Fix a bounded window
$K\subset\overline{\HH}$ and buffers $\rho_-,\rho_+>0$.  Deterministically partition $J$
into $O(h^{-1})$ intervals of length at most $h$, with endpoints off the
graph.  Let $E_{\delta,h}$ be the event that, for one of these intervals
$I$, one directed half-plane cycle has two portions $\alpha,\beta$
contained in $K$ such that each portion has one endpoint in
$\{u\le c-\rho_-\}$ and the other in $\{u\ge c+\rho_+\}$, and
\[
 \iota_I(\alpha)\iota_I(\beta)<0.
\]
Then
\begin{equation}
 \lim_{h\downarrow0}\limsup_{\delta\downarrow0}
 \PP(E_{\delta,h})=0.
 \label{eq:local-opposite-limit}
\end{equation}
For a vertical line $L$, the segment may reach or extend below the
bottom boundary.  Test lines are chosen off lattice vertices; the
same estimate holds for seam perturbations of order $\delta$.
\end{proposition}

\begin{proof}
Fix a coupling error $\eps>0$ and apply
Proposition~\ref{prop:localization} to a neighborhood of $K$.
On its agreement event, both directed portions remain rectangle
paths with exactly the same signs, because the two matching labels
are preserved.  There are two possibilities for their new connections.
If the portions belong to distinct rectangle cycles, their two nonzero
indices give the second forbidden event in
Proposition~\ref{prop:rectangle-short}.  If they belong to the same
rectangle cycle $g$, its whole-cycle index satisfies
$\iota_I(g)\in\{-1,0,1\}$ by the Jordan curve theorem.  Two opposite nonzero
integers cannot both equal $\iota_I(g)$; at least one portion therefore
gives the first forbidden event of that proposition.

The rectangle can be chosen with all required cuts longitudinally
interior.  For a horizontal middle line, a nonempty event already
requires endpoints below it; take its lower cap between those
endpoints and the middle line, at positive height.  For a vertical
middle line there is no corresponding restriction at the bottom:
the test interval is transverse, and
Proposition~\ref{prop:rectangle-short} permits its truncation at a
transverse end.  Rotation and reflection give the same argument in
both coordinate directions and for either direction of each portion.
Consequently
\[
 \lim_{h\downarrow0}\limsup_{\delta\downarrow0}
 \PP(E_{\delta,h})\le\eps.
\]
Let $\eps\downarrow0$.  Rounding a cut by $O(\delta)$ leaves fixed
positive buffers and longitudinal room, and the rectangle bounds are
uniform under this rounding.  Alternatively, along any prescribed
mesh sequence one may choose fixed test lines and endpoints avoiding
all its lattice sites, as we shall do in Section~\ref{sec:identification}.
\end{proof}

\subsection{From local traversals to compactness in the disk}

We now use Corollary~\ref{cor:local-crossings} to control parametrizations,
not merely the number of loops. The use of crossing control and
multiscale time changes to obtain compactness follows the strategy of
Aizenman--Burchard~\cite[Lemma~2.4 and Section~4]{AB}. We use the direct
partition construction below, which requires only our fixed-scale
traversal tightness. For a collection of parametrized loops
in $\overline{\DD}$, let $N_\eta$ be the largest total number of
parameter-disjoint portions whose two endpoints have Euclidean
distance greater than $\eta$.  Parameter interiors are disjoint along each individual circle; common
endpoints are allowed. Disjointness is automatic between different loops.
This number is invariant under reparametrization and reversal.

\begin{lemma}[Detection of disk trips]
\label{lem:disk-trips}
For every $\eta>0$, the random variables
$N_\eta(F(\mathcal L_\delta))$ are tight as $\delta\downarrow0$.
\end{lemma}

\begin{proof}
The formulas
\begin{equation}
 |F(z)-1|=\frac{2}{|z+i|},
 \qquad |F'(z)|\le2\quad(z\in\overline{\HH})
 \label{eq:compactification-control}
\end{equation}
give a fixed bounded window in which every macroscopic trip can be
detected.  Indeed, if two disk endpoints are more than $\eta$ apart,
one is more than $\eta/2$ from $1$.  Start the portion at that endpoint,
reversing it if necessary, and stop when its disk displacement first
reaches $\eta/4$.  Until then it stays at distance greater than
$\eta/4$ from $1$, hence its inverse image lies in a bounded half-plane
window depending only on $\eta$.  Its Euclidean endpoint displacement
in the half-plane is at least $\eta/8$ by
\eqref{eq:compactification-control}.  One coordinate therefore changes
by at least $\eta/(8\sqrt2)$.

Choose a rectangular grid with spacing much smaller than this last
quantity in a slightly larger fixed window.  The portion must traverse
between a pair of separated grid lines in one coordinate direction.
There are only finitely many possible slabs.  For a vertical
displacement the two detecting horizontal lines can both be taken at
positive height, even when the portion starts near the bottom.
Horizontal displacement is detected by vertical lines, whose
transverse range may include the bottom.  Perturb the finitely many
lines to seams when needed.

Restricting parameter-disjoint trips to these detecting portions
preserves disjointness.  Assign each one a slab it traverses; their
number is bounded by the sum of the finitely many local traversal
counts in Corollary~\ref{cor:local-crossings}.  This proves tightness.
\end{proof}

\begin{proposition}[Preliminary compactness]
\label{prop:tightness}
The laws of $F(\mathcal L_\delta)$ are tight in the collection space of
closed-disk curves with the ensemble matching topology, with constant
loops discarded.  Moreover, the loops can be ordered by decreasing
diameter and given measurable, orientation-preserving parametrizations
whose joint laws are tight in the product of uniform path spaces.
Every subsequential limit has finitely many loops above every positive
diameter, and convergence retains all its nonconstant entries.
At this stage limiting curves may be nonsimple, touch the boundary,
or occur with multiplicity.
\end{proposition}

\begin{proof}
We will choose measurable parametrized representatives for the joint
limit in Section~\ref{sec:identification} and also prove compactness in
the collection metric. The same trip bounds will give both conclusions.
First, every loop of diameter greater than $\eta$ contributes a trip
at that scale, so its number is at most $N_\eta$.
For a fixed mesh these numbers, and all the trip counts, are finite:
the detection argument forces each such trip to use a portion in a
bounded lattice window, and the cycles use each lattice edge at most
once.  In fact, for $\delta$ bounded away from zero and bounded above,
these counts are bounded deterministically.  The total length of
lattice edges meeting the detecting window is uniformly bounded in
that range of meshes, whereas each detecting portion has a fixed
positive endpoint displacement.

Choose the parametrization of every loop by the following single rule.
Start with the image under $F$
of its usual polygonal parametrization,
$\gamma:[0,1]\to\overline{\DD}$, rooted by a fixed measurable rule and
oriented by its matching labels.  At scale
$r_j=2^{-j}$, partition its parameter interval greedily: from the last
partition point stop at the first displacement $r_j$, repeating until
time $1$.  Delete an empty final interval if necessary.  The resulting
partition $\mathcal P_j$ has
\[
 q_j:=|\mathcal P_j|\le N_{r_j/2}+1,
\]
and each piece has image diameter at most $2r_j$.
Let $w_j=2^{-j-1}$ and put on $[0,1]$ the probability measure
\begin{equation}
 d\nu_\gamma(t)=\frac12\,dt+
 \sum_{j\ge1}\frac{w_j}{q_j}
       \sum_{I\in\mathcal P_j}\frac{\1_I(t)}{|I|}\,dt.
 \label{eq:partition-mass}
\end{equation}
The cumulative function $\theta_\gamma(t)=\nu_\gamma([0,t])$ is a
strictly increasing continuous homeomorphism.  Use
$\widetilde\gamma=\gamma\circ\theta_\gamma^{-1}$ as the new
parametrization. Greedy stopping times, the measures in
\eqref{eq:partition-mass}, and their inverses are measurable functions
of the original finite polygonal paths. This rule depends only on the
loop, not on any probability tolerance or trip-count bound.

Now fix $\eps>0$. Lemma~\ref{lem:disk-trips}, together with the
deterministic bounds for meshes bounded away from zero, lets us choose
integers $A_j<\infty$ such that
\begin{equation}
 \PP\bigl(N_{2^{-j-1}}\le A_j\text{ for every }j\ge1\bigr)
 \ge1-\eps
 \label{eq:all-trip-scales}
\end{equation}
uniformly for $0<\delta\le1$. Allocate error $\eps2^{-j}$ at level
$j$ and use a union bound. On this event, $q_j\le A_j+1$ for every
loop and every $j$.

Each interval of $\mathcal P_j$ now has duration at least
$w_j/q_j\ge w_j/(A_j+1)$.  A time interval shorter than this cannot
contain a full partition piece and hence meets at most two adjacent
pieces.  Its image has diameter at most $4r_j$.  Thus
\eqref{eq:all-trip-scales} gives a common modulus of continuity for
every reparametrized loop.  The same argument applies across the
identified endpoints of the parameter circle.

Arzel\`a--Ascoli now gives a compact set of parametrized representatives
on the event \eqref{eq:all-trip-scales}.  Order the loops by decreasing
diameter, resolving ties by a fixed enumeration of the finite lattice
polygons rescaled to mesh one, and pad a finite list by constant curves.
This is measurable:
there are only finitely many loops above every positive diameter, so
successive largest entries exist and exhaust the list.  The joint
ordered representatives lie in a compact product, and the diameter
counts give a uniform vanishing bound on the tails of the lists.

For completeness, consider any sequence of lists with these bounds.
Take a diagonal subsequence on which each parametrized entry converges
uniformly.  The limit list has only finitely many entries of diameter
greater than any prescribed positive number, by using the count bound
at a smaller cutoff.  Discard its constant entries.  At a given error
scale, a fixed finite initial list contains all larger loops on both
sides.  Match its corresponding nonconstant entries; entries with
constant limit are eventually smaller than the error scale.  Uniform
convergence on this finite list and the vanishing diameter tails give
the required partial matching.  This proves relative compactness in
the ensemble topology as well as tightness of the parametrized lists.
The compact closures supplied by these bounds have probability at
least $1-\eps$, proving the proposition.
\end{proof}

The distinction between this proposition and the desired conclusion is
essential.  Traversal bounds retain every macroscopic path, but they
do not by themselves exclude collapsed interiors or retracing.
Proposition~\ref{prop:local-opposite}, together with the fixed-puncture
law, will rule out those possibilities and identify every surviving
curve in Section~\ref{sec:identification}.
